% Part: lambda-calculus
% Chapter: introduction
% Section: lambda-definability

\documentclass[../../../include/open-logic-section]{subfiles}

\begin{document}

\olfileid{lam}{int}{rep}
\olsection{\usetoken{S}{lambda definable} পাটিগাণিতিক অপেক্ষক}

ল্যাম্বডা ক্যালকুলাস কীভাবে গণনার একটি মডেল হিসেবে কাজ করতে পারে? প্রথমে
এই বক্তব্যটির অর্থ কী হবে, তা-ই স্পষ্ট নয়। স্বাভাবিক সংখ্যার উপর
গণনসাধ্যতা নিয়ে কথা বলতে হলে এই সংখ্যাগুলির একটি উপযুক্ত উপস্থাপন খুঁজে
নিতে হবে। নিচের উপস্থাপনটি আশ্চর্যজনকভাবে ভালো কাজ করে।

\begin{defn}
প্রতিটি স্বাভাবিক সংখ্যা~$n$-এর জন্য \emph{চার্চ সংখ্যাপদ}
$\num{n}$-কে ল্যাম্বডা পদ $\lambd[x][\lambd[y][(x(x(x(\dots
x(y)))))]]$ হিসেবে সংজ্ঞায়িত করো, যেখানে মোট $n$টি~$x$ আছে।
\end{defn}

$\num{n}$ পদগুলি ``পুনরাবর্তক'': নিবেশ~$f$-এর জন্য $\num{n}$ এমন অপেক্ষক
ফেরত দেয়, যা $y$-কে~$f^n(y)$-তে পাঠায়। লক্ষ করো, প্রতিটি সংখ্যাপদই
স্বাভাবিক রূপে আছে। এবার স্বাভাবিক সংখ্যার উপর কোনো অপেক্ষককে একটি
ল্যাম্বডা পদ ``গণনা করে'' বলতে কী বোঝায়, তা বলা যায়।

\begin{defn}
ধরা যাক $f(x_0, \dots, x_{k-1})$ একটি $k$-স্থানীয় আংশিক অপেক্ষক, যা
$\Nat$ থেকে $\Nat$-এ যায়। বলি, কোনো $\lambd$-পদ~$F$
অপেক্ষক~$f$-কে \emph{!!{lambda define}s}, যদি স্বাভাবিক সংখ্যার প্রতিটি
ক্রম $n_0$, \dots,~$n_{k-1}$-এর জন্য
\[
F\, \num{n_0}\, \num{n_1} \dots \num{n_{k-1}} \red \num{f(n_0, n_1, \dots,
  n_{k-1})}
\]
যখন $f(n_0, \dots, n_{k-1})$ সংজ্ঞায়িত, এবং অন্যথায় $F\, \num{n_0}\,
\num{n_1} \dots \num{n_{k-1}}$-এর কোনো স্বাভাবিক রূপ নেই।
% BN-SRC-268 TARGET-CORRECTION-BEGIN
\par\smallskip\noindent\textnormal{\small\textbf{উৎস-সংশোধন (BN-SRC-268):} হিমায়িত সংজ্ঞায় চলরাশি ও নিবেশগুলি~$0$ থেকে~$k-1$ পর্যন্ত সূচিত হলেও অপেক্ষকটিকে অসঙ্গতভাবে~$n$-স্থানীয় বলা হয়েছে। বাংলা পাঠে স্থানসংখ্যা~$k$ পুনরুদ্ধার করা হয়েছে।} % BN-SRC-268 TARGET-CORRECTION-END
% BN-SRC-269 TARGET-CORRECTION-BEGIN
\par\smallskip\noindent\textnormal{\small\textbf{উৎস-সংশোধন (BN-SRC-269):} অসংজ্ঞায়িত অবস্থার পদটিতে হিমায়িত সূত্র~$F$ ও প্রথম সংখ্যাপদের মধ্যে প্রয়োগের বদলে কমা বসিয়েছে। বাংলা সূত্রে কমাটি সরিয়ে প্রয়োগের ফাঁক পুনরুদ্ধার করা হয়েছে, যাতে একই~$F\,\num{n_0}\dots\num{n_{k-1}}$ পদটির স্বাভাবিক রূপ না থাকার কথা বলা হয়।} % BN-SRC-269 TARGET-CORRECTION-END
\end{defn}

\begin{thm}
\ollabel{thm:lambda-def}
কোনো অপেক্ষক~$f$ আংশিক গণনসাধ্য যদি এবং কেবল যদি কোনো ল্যাম্বডা পদ
তাকে !!{lambda defined} করে।
\end{thm}

\begin{explain}
এই উপপাদ্যটি বেশ চমকপ্রদ। গণনার মডেল হিসেবে ল্যাম্বডা ক্যালকুলাস খুবই
সরল একটি ক্যালকুলাস; এর একমাত্র ক্রিয়া হলো ল্যাম্বডা-বিমূর্তন ও প্রয়োগ!{}
তবুও এই সামান্য উপকরণ দিয়েই যেকোনো গণনপদ্ধতি বাস্তবায়ন করা সম্ভব।
\end{explain}

\end{document}
